% Solution: Tadelis (2013) Exercise 5.1 "Prove Proposition 5.1". Input by \solitem. Check: check_L2.py.
\begin{exsol}{Tadelis Exercise 5.1}{95}{Exam}{3}
\asked Prove Proposition~5.1: if a strategy profile $s$ is a strict dominant-strategy equilibrium, or
the unique survivor of IESDS, or the unique rationalizable profile, then $s$ is the unique Nash
equilibrium.

\method Direct arguments. For (1) and (2), a profitable deviation against $s_{-i}$ forces the
deviating strategy to have been deleted at some round; for (3), the definition of rationalizability
applied to a singleton. The slides prove (2) in two halves, Propositions~1 and~2 (Lecture~2,
slides~16--17), and the proof below follows the same ``first round in which something is deleted''
argument. All games are finite; ``Nash equilibrium'' means a pure-strategy one, as in Tadelis \S5.1.

\solution
\textbf{Part 1: strict dominant-strategy equilibrium.} Let $s_i^*$ be strictly dominant for every
$i$: $v_i(s_i^*, s_{-i}) > v_i(s_i, s_{-i})$ for all $s_i \ne s_i^*$ and all $s_{-i}$.
\begin{itemize}
  \item \emph{$s^*$ is a Nash equilibrium.} Take $s_{-i} = s^*_{-i}$: the inequality says $s_i^*$
  does strictly better than any other $s_i$, so it is a best reply to $s^*_{-i}$ for every $i$.
  \item \emph{It is the only one.} If $s \ne s^*$ were a Nash equilibrium, some $s_i \ne s_i^*$ would
  be a best reply to $s_{-i}$, contradicting $v_i(s_i^*, s_{-i}) > v_i(s_i, s_{-i})$.
\end{itemize}
(Part 1 also follows from Part 2: a strict dominant strategy deletes all the others in round~1.)

\textbf{Part 2: unique survivor of IESDS.} Let $S^k$ be the strategy sets after $k$ rounds and let
$s^*$ be the only profile in $S^\infty$.
\begin{itemize}
  \item \emph{$s^*$ is a Nash equilibrium (Proposition~2).} Suppose player $i$ had a strategy $s_i$
  with $v_i(s_i, s^*_{-i}) > v_i(s^*_i, s^*_{-i})$. Then $s_i \ne s_i^*$, so $s_i$ was deleted in
  some round $k$: there is $s_i' \in S_i^{k-1}$ with $v_i(s_i', t_{-i}) > v_i(s_i, t_{-i})$ for all
  $t_{-i} \in S_{-i}^{k-1}$. The profile $s^*_{-i}$ lies in $S_{-i}^{k-1}$ (it survives forever), so
  \[
    v_i(s_i', s^*_{-i}) > v_i(s_i, s^*_{-i}) > v_i(s_i^*, s^*_{-i}),
  \]
  and $s_i' \ne s_i^*$ is again a strictly better reply to $s^*_{-i}$. Apply the same argument to
  $s_i'$, and so on. The values $v_i(\cdot, s^*_{-i})$ rise strictly along the chain, so no strategy
  repeats, and the strategy sets are finite: the chain must stop at a strategy that is never
  deleted, that is, at $s_i^*$. But every member of the chain does strictly better against
  $s^*_{-i}$ than $s_i^*$ does, a contradiction. So no profitable deviation exists.
  \item \emph{It is the only Nash equilibrium (Proposition~1).} Let $s$ be a Nash equilibrium and
  suppose it does not survive. Take the first round $k$ in which some component $s_i$ is deleted:
  $s_i$ is strictly dominated by some $s_i' \in S_i^{k-1}$ against every $t_{-i} \in S_{-i}^{k-1}$,
  and $s_{-i} \in S_{-i}^{k-1}$ because no component was deleted earlier. Then
  $v_i(s_i', s_{-i}) > v_i(s_i, s_{-i})$, so $s_i$ is not a best reply to $s_{-i}$, a contradiction.
  Hence every Nash equilibrium lies in $S^\infty = \{s^*\}$.
\end{itemize}

\textbf{Part 3: unique rationalizable profile.} Let $R^k$ be the strategy sets after $k$ rounds of
deleting strategies that are never a best reply to a belief over the opponents' surviving
strategies, and let $R^\infty = \{s^*\}$.
\begin{itemize}
  \item \emph{$s^*$ is a Nash equilibrium.} The sets are finite and decrease to $R^\infty$, so for
  $k$ large $R^{k-1}_{-i} = \{s^*_{-i}\}$. Since $s_i^*$ survives round $k$, it is a best reply to some
  belief over $R^{k-1}_{-i}$, and the only such belief is the point mass on $s^*_{-i}$. So $s_i^*$ is a
  best reply to $s^*_{-i}$, for each $i$.
  \item \emph{It is the only one: every Nash equilibrium is rationalizable.} Let $s$ be a Nash
  equilibrium. By induction on $k$, $s \in R^k$: it holds for $k = 0$; if $s_{-i} \in R^{k-1}_{-i}$,
  then $s_i$ is a best reply to the belief ``$s_{-i}$ with probability one'', which is a belief over
  $R^{k-1}_{-i}$, so $s_i \in R^k_i$. Hence $s \in R^\infty = \{s^*\}$.
\end{itemize}
$\blacksquare$

\interp The three hypotheses are three ways of saying that rational, mutually aware players have a
unique prediction, and they are nested: a strict dominant-strategy equilibrium is the unique
survivor of IESDS (everything else goes in round~1), and a unique IESDS survivor is the unique
rationalizable profile (MWG exercise 8.C.1). So Nash equilibrium is a \emph{necessary} condition for
any of them to make a single prediction; this is why the slides treat IESDS and rationalizability as
the epistemic benchmark and Nash equilibrium as the sharper concept (Lecture~2, slides~15--17). The
converse fails: the Battle of the Sexes has two Nash equilibria and IESDS deletes nothing.

\verify \texttt{check\_L2.py} tests the proposition on random finite two-player games ($1{,}200$
games with $2$--$4$ strategies per player and payoffs in $0,\dots,5$, ties included): whenever the
game has a strict dominant-strategy equilibrium, a unique IESDS survivor (pure dominators), or a
unique rationalizable profile (computed by linear programming over mixed beliefs), the set of pure
Nash equilibria is exactly that profile. It also confirms the lemma that every pure Nash equilibrium
is rationalizable. The proof itself is the argument above. Tadelis gives no answer.

\mistakes
\begin{itemize}
  \item Proving only ``$s$ is a Nash equilibrium'' and forgetting uniqueness, which is the other
  half (Proposition~1).
  \item Using \emph{weak} dominance: a weakly dominant-strategy equilibrium can coexist with other
  Nash equilibria (exercise 5.2).
  \item Assuming that the strategy deleted in some round is dominated by $s_i^*$ itself; the dominator
  may be deleted later, which is why the chain argument is needed.
\end{itemize}
\end{exsol}
